\section{Результаты}
\label{sec:results}

Для каждого из 14 компонентов проведены четыре серии сравнений
(\texttt{cpu\_vs\_xla}, \texttt{xla\_vs\_cublas}, \texttt{xla\_vs\_cuda},
\texttt{cublas\_vs\_cuda}). GPU-реализации исполнялись на NVIDIA A100; время
компиляции графа в замеры не включалось.

\subsection{Ускорение симуляции}

Ускорение определяется долей GEMM в графе. Трансформерный блок ускоряется вдвое,
полносвязная сеть~--- вчетверо (табл.~\ref{tab:speedup}); cuBLAS и CUDA-ядро по
скорости равны ($<1\%$). Явное управление порядком суммирования в CUDA-ядре не
снижает производительности.

\begin{table}[htbp]
\centering
\caption{Ускорение симуляции при переносе GEMM на GPU.}
\label{tab:speedup}
\fitwidth{\begin{tabular}{l c c c c c}
    \hline
    & \multicolumn{3}{c}{\textbf{Время симуляции, с}}
    & \multicolumn{2}{c}{\textbf{Ускорение к CPU}} \\
    \cmidrule(lr){2-4}\cmidrule(lr){5-6}
    \textbf{Модель} & \textbf{CPU} & \textbf{cuBLAS} & \textbf{CUDA} &
    \textbf{cuBLAS} & \textbf{CUDA} \\
    \hline
    Трансформер (2 слоя)        & 23.03 & 11.93 & 12.07 & $1.93\times$ & $1.91\times$ \\
    Полносвязная сеть (6 слоёв) & 14.70 & 3.84  & 3.89  & $3.83\times$ & $3.78\times$ \\
    \hline
\end{tabular}}
\end{table}

\subsection{Послойный профиль и разделение вкладов}

Во всех 14 компонентах профиль однороден и трёхуровневый. GEMM-слои
(\texttt{nn.Linear}, \texttt{nn.Conv2d}, проекции Q/K/V, $AV$, FFN) дают
$d_{\text{rel}}^{(l)}\approx2.2\cdot10^{-3}$ (диапазон по всем слоям
$2.17$--$2.37\cdot10^{-3}$, вариация $<10\%$); нормализационные слои~---
$8$--$9\cdot10^{-8}$ (разрыв в \emph{четыре порядка}); точные операции (RoPE,
тиражирование KV, Dropout, Flatten)~--- ровно $0$.

Послойную структуру наглядно демонстрирует RoPE Attention
(табл.~\ref{tab:rope}): проекции Q/K/V, взвешенное суммирование $AV$ и слои
FFN дают ${\approx}2.2$--$2.3\cdot10^{-3}$; RMSNorm~--- ${\approx}8$--$9\cdot10^{-8}$;
поворотное кодирование и Dropout (с фиксированным генератором)~--- ровно $0$, так
как не выполняют арифметики над данными. Отдельного комментария требует шаг
вычисления каузальной матрицы оценок внимания: его $d_{\text{rel}}\approx0$ в
прямом проходе~--- артефакт метрики, а не физический результат. Каузальная маска
заполняет ${\approx}50\%$ выхода значением $-10^9$, раздувая знаменатель
$\|\Vec{y}^{(l)}\|_F$ при нулевом вкладе замаскированных позиций в числитель; по
незамаскированным позициям $d_{\text{rel}}\approx2.3\cdot10^{-3}$, что
совпадает с уровнем GEMM (его и показывает свободный от маскирования обратный
проход).

\begin{table}[htbp]
\centering
\caption{Послойные расхождения RoPE Attention (\texttt{cpu\_vs\_xla}), фрагмент;
    столбцы~--- $d_{\text{rel}}$ в прямом~(пр.) и обратном~(обр.) проходе.}
\label{tab:rope}
\fitwidth{\begin{tabular}{l l c c}
    \hline
    \textbf{Модуль} & \textbf{Тип} &
    \textbf{$d_{\text{rel}}$ (пр.)} & \textbf{$d_{\text{rel}}$ (обр.)} \\
    \hline
    RMSNorm             & Norm        & $7.98\cdot10^{-8}$ & $1.05\cdot10^{-7}$ \\
    Linear (Q, K, V)    & GEMM        & ${\approx}2.21\cdot10^{-3}$ & ${\approx}2.22\cdot10^{-3}$ \\
    ApplyRoPE           & Exact       & $0$ & $0$ \\
    CausalAttnScores    & AttnScores  & ${\approx}0^{*}$ & $2.339\cdot10^{-3}$ \\
    WeightedValue ($AV$) & GEMM       & $2.331\cdot10^{-3}$ & $2.337\cdot10^{-3}$ \\
    Linear (out)        & GEMM        & $2.250\cdot10^{-3}$ & $2.222\cdot10^{-3}$ \\
    FFN up / down       & GEMM        & ${\approx}2.26\cdot10^{-3}$ & ${\approx}2.28\cdot10^{-3}$ \\
    Dropout             & Exact       & $0$ & $0$ \\
    \hline
    \multicolumn{4}{l}{\footnotesize $^{*}$~артефакт каузальной маски; по незамаскированным позициям ${\approx}2.3\cdot10^{-3}$.} \\
\end{tabular}}
\end{table}

Уровень $2.2\cdot10^{-3}$ формируется не при умножении матриц, а при компиляции
графа HLO (слияние операций, переназначение раскладки памяти, переупорядочивание
$K$-редукции). Поскольку граф HLO одинаков для всех трёх реализаций,
попарное сравнение XLA-реализаций изолирует вклад GEMM
(табл.~\ref{tab:levels}): он на 2--5 порядков ниже компиляторного.

\begin{table}[htbp]
\centering
\caption{Три уровня расхождения GEMM-слоёв (прямой проход).}
\label{tab:levels}
\fitwidth{\begin{tabular}{l c l}
    \hline
    \textbf{Пара сравнения} & \textbf{$d_{\text{rel}}$} & \textbf{Источник} \\
    \hline
    \texttt{cpu} vs \texttt{xla\_*}        & ${\sim}2.2\cdot10^{-3}$ & Компилятор XLA (весь стек) \\
    \texttt{xla\_cpu} vs \texttt{xla\_cuda}   & ${\sim}1.2\cdot10^{-5}$ & Реализация: CUDA-ядро \\
    \texttt{xla\_cublas} vs \texttt{xla\_cuda} & ${\sim}1.2\cdot10^{-5}$ & Две GPU-реализации между собой \\
    \texttt{xla\_cpu} vs \texttt{xla\_cublas} & ${\sim}9\cdot10^{-8}$     & Реализация: cuBLAS \\
    \hline
\end{tabular}}
\end{table}

\begin{figure}[H]
\centering
\resizebox{0.92\textwidth}{!}{%
\begin{tikzpicture}[font=\small]
    % ось: десятичный логарифм d_rel, 2.2 см на декаду, -8 -> 0
    \draw[-{Stealth[length=5pt]}, line width=0.8pt] (-0.3,0) -- (12.7,0)
        node[right, font=\footnotesize] {$d_{\text{rel}}$};
    \foreach \e/\x in {{-8}/0, {-7}/2.2, {-6}/4.4, {-5}/6.6, {-4}/8.8, {-3}/11.0} {
        \draw[line width=0.6pt] (\x,0.08) -- (\x,-0.08) node[below, font=\scriptsize] {$10^{\e}$};
    }
    % epsilon_FP32 = 1.2e-7 -> x = 2.37
    \draw[dashed, gray, line width=0.6pt] (2.37,-0.15) -- (2.37,1.55);
    \node[above, font=\scriptsize, gray!55!black] at (2.37,1.55) {$\varepsilon_{\texttt{FP32}}$};
    % cuBLAS 9e-8 -> x=2.10
    \draw[green!55!black, line width=1pt] (2.10,0) -- (2.10,0.85);
    \fill[green!55!black] (2.10,0) circle (2.6pt);
    \node[above, font=\scriptsize, green!45!black, align=center] at (2.10,0.85) {cuBLAS\\[-1pt]$9{\cdot}10^{-8}$};
    % CUDA 1.2e-5 -> x=6.77
    \draw[orange!80!black, line width=1pt] (6.77,0) -- (6.77,1.0);
    \fill[orange!80!black] (6.77,0) circle (2.6pt);
    \node[above, font=\scriptsize, orange!60!black, align=center] at (6.77,1.0) {CUDA-ядро\\[-1pt]$1.2{\cdot}10^{-5}$};
    % Компилятор 2.2e-3 -> x=11.75
    \draw[red!75!black, line width=1pt] (11.75,0) -- (11.75,1.3);
    \fill[red!75!black] (11.75,0) circle (2.6pt);
    \node[above, font=\scriptsize, red!60!black, align=center] at (11.4,1.3) {\textbf{Компилятор XLA}\\[-1pt]$2.2{\cdot}10^{-3}$};
    % брэйс: межреализационный диапазон GEMM
    \draw[decorate, decoration={brace, amplitude=4pt}, line width=0.6pt]
        (2.10,-0.62) -- (6.77,-0.62)
        node[midway, below, font=\scriptsize, yshift=-1pt, align=center] {различия реализаций GEMM};
\end{tikzpicture}}
\caption{Иерархия источников расхождения GEMM-слоёв по величине послойного
    отклонения $d_{\text{rel}}$ (логарифмическая шкала). Вклад компилятора XLA на
    2--5 порядков превосходит различия между реализациями GEMM; отклонение cuBLAS
    лежит у машинного эпсилона \texttt{FP32}, поэтому замена CPU-реализации на GPU
    допустима по точности.}
\label{fig:hierarchy}
\end{figure}

\textbf{Ключевой вывод:} доминирует компилятор (рис.~\ref{fig:hierarchy}).
Замена CPU-GEMM на cuBLAS или CUDA-ядро при той же цепочке компиляции меняет тип
ошибки, но не её доминирующую компоненту. Вопреки ожиданиям, cuBLAS оказывается на два порядка ближе к CPU, чем
CUDA-ядро (более консервативный порядок редукций; отклонение cuBLAS~--- порядка
одной единицы последнего разряда \texttt{FP32}). Слабая зависимость
$d_{\text{rel}}$ от $K$ (отношение $1.05$ при $K_2/K_1=3.5$ вместо
$\sqrt{3.5}\approx1.87$) согласуется с доминированием компилятора над
GEMM-арифметикой.

\subsection{Диагностика природы ошибки}

Метрика смещения $b$~\eqref{eq:diag} однозначно разделяет пары
(табл.~\ref{tab:diag}). Компиляторная ошибка (\texttt{cpu\_vs\_xla}) случайна и
аддитивна ($b\approx0.5$): GEMM суммирует $\sim K$ независимых ошибок округления
с нулевым средним. Единственная смещённая пара~--- \texttt{xla\_vs\_cublas}
($b\approx0.28$, более $70\%$ элементов одного знака): фиксированный порядок
обхода операндов в cuBLAS создаёт направленный сдвиг относительно CPU-GEMM.
CUDA-ядро смещения не создаёт ни относительно CPU, ни относительно cuBLAS.

\begin{table}[htbp]
\centering
\caption{Диагностика природы расхождений GEMM-слоёв (прямой проход).}
\label{tab:diag}
\fitwidth{\begin{tabular}{l c c l}
    \hline
    \textbf{Пара} & \textbf{$b$} & \textbf{$c$} & \textbf{Классификация} \\
    \hline
    \texttt{cpu\_vs\_xla}     & $0.498$--$0.522$ & $-0.06$--$0.10$ & Случайная аддитивная \\
    \texttt{xla\_vs\_cublas}  & $0.27$--$0.34$   & $0.08$--$0.16$  & \textbf{Смещённая} аддитивная \\
    \texttt{xla\_vs\_cuda}    & $0.49$--$0.51$   & $-0.06$--$0.10$ & Случайная аддитивная \\
    \texttt{cublas\_vs\_cuda} & $0.49$--$0.51$   & $-0.06$--$0.10$ & Случайная аддитивная \\
    \hline
\end{tabular}}
\end{table}

\subsection{Накопление: прямой проход}

Малое послойное отклонение необходимо, но не достаточно для приемлемости
реализации: даже небольшая ошибка способна накопиться до значимого масштаба при
прохождении многослойной архитектуры, причём характер накопления зависит от
архитектуры. Поэтому требуется явный анализ накопления. Коэффициент
$R_{\text{fwd}}$~\eqref{eq:accum} образует группы по типу нормализации
(табл.~\ref{tab:rfwd}). Pre-Norm LayerNorm (RoBERTa) даёт $R=0.53<1$:
центрирование $(I-\tfrac{1}{d}\Vec{1}\Vec{1}^\top)$ обнуляет равномерную
составляющую расхождения на входе каждого подслоя. RMSNorm среднее не вычитает,
поэтому равномерная составляющая проходит и $R=1.4$--$1.7$; максимум у FFN
SwiGLU/GeGLU (нет сжимающего softmax-внимания), минимум у Post-LN DeepNorm
(нормализация после суммирующей связи $\Vec{y}=\Vec{x}+F(\Vec{x})$~--- остаточного
соединения, по которому расхождение подслоя $F$ добавляется к входу напрямую,
минуя нормировку). LeNet без нормализации удерживается вблизи $R\approx1$ за счёт
Sigmoid-стоков.

\begin{table}[htbp]
\centering
\caption{Коэффициент накопления $R_{\text{fwd}}$ по группам компонентов
    ($^\dagger$~для \texttt{xla\_vs\_cublas} приведены оценки порядка величины:
    сигнал на уровне $\varepsilon_{\texttt{FP32}}$).}
\label{tab:rfwd}
\fitwidth{\begin{tabular}{l l c c c}
    \hline
    \textbf{Компонент / группа} & \textbf{Нормализация} &
    \textbf{\texttt{cpu\_vs\_xla}} & \textbf{\texttt{xla\_vs\_cuda}} &
    \textbf{\texttt{xla\_vs\_cublas}$^\dagger$} \\
    \hline
    RoBERTa                & Pre-Norm LN  & $\mathbf{0.53}$ & $0.43$ & ${\sim}2.0$ \\
    LeNet                  & нет          & $1.08$ & $0.81$ & ${\sim}1.3$ \\
    Внимание (8 шт.)       & RMSNorm      & $1.42$--$1.46$ & $1.09$--$1.30$ & ${\sim}4.5$--$4.7$ \\
    RMSNorm FFN            & RMSNorm      & $1.47$ & $1.32$ & ${\sim}4.7$ \\
    SwiGLU / GeGLU FFN     & RMSNorm      & $1.70$--$1.71$ & $1.48$--$1.51$ & ${\sim}4.1$ \\
    DeepNorm               & Post-LN      & $1.39$ & $1.15$ & ${\sim}4.2$ \\
    \hline
\end{tabular}}
\end{table}

Для случайной ошибки (\texttt{cpu\_vs\_xla}, \texttt{xla\_vs\_cuda}) накопление
умеренное: $R_{\text{fwd}}\le1.7$ во всех архитектурах. Метрика $b$ оказывается
более информативным предиктором накопления, чем норма $d_{\text{rel}}$:
\texttt{xla\_vs\_cublas} ($b\approx0.28$, $d_{\text{rel}}\sim10^{-8}$) формально
даёт $R\approx4$, тогда как \texttt{xla\_vs\_cuda} ($b\approx0.5$,
$d_{\text{rel}}\sim10^{-5}$)~--- лишь $R\approx1.4$. Однако послойный сигнал
cuBLAS лежит на уровне $\varepsilon_{\texttt{FP32}}$, поэтому значения
$R_{\text{fwd}}$ для этой пары трактуются лишь как оценки порядка величины.

\subsection{Накопление: обратный проход}

В изолированном режиме («обратный») накопление умеренное и сопоставимо с прямым
проходом: $R_{\text{bwd}}\approx1.4$--$1.9$ для случайной ошибки
(табл.~\ref{tab:rbwd}). Исключение~--- DeepNorm ($R_{\text{bwd}}=0.35$): Post-LN
выступает стоком градиентных расхождений симметрично своей роли в прямом проходе.
В совместном режиме присутствуют обе составляющие
разложения~\eqref{eq:grad-decomp}: отклонённые промежуточные тензоры прямого
прохода изменяют якобианы при вычислении градиентов, добавляя
$\Delta_{\text{fwd}}$ к собственной ошибке обратного прохода
$\Delta_{\text{bwd}}$. Поскольку эти составляющие не компенсируют друг друга,
накопление возрастает до $R_{\text{bwd}}\approx5$--$6$~--- втрое выше
изолированного режима. Это подтверждает, что усиление обусловлено совместным
действием отклонений обоих проходов, а не спецификой градиентных вычислений.
Аномалия RMSNorm~FFN ($R_{\text{bwd}}=17.02$) вызвана отсутствием суммирующей
связи $\Vec{y}=\Vec{x}+F(\Vec{x})$: отклонённый вход нелинейно меняет множитель
$1/\mathrm{RMS}(\Vec{x})$, проходящий по градиентной цепочке без ослабления;
в компонентах с суммирующей связью прямая ветвь $\partial\mathcal{L}/\partial\Vec{y}$
минует якобиан подслоя $J_F$ и частично ослабляет суммарное отклонение. Несмотря
на экстремальное значение, линейное приближение остаётся применимым: для
доминирующих GEMM-слоёв обратный проход $\nabla_{\Vec{x}}=\nabla_{\Vec{y}}W^\top$
не зависит от выходов прямого прохода, и накопление обусловлено аддитивным
суммированием погрешностей, а не нелинейным взаимодействием.

\begin{table}[htbp]
\centering
\caption{Коэффициент накопления $R_{\text{bwd}}$ по двум режимам
    (\texttt{cpu\_vs\_xla}).}
\label{tab:rbwd}
\fitwidth{\begin{tabular}{l c c}
    \hline
    \textbf{Компонент / группа} & \textbf{Обратный} & \textbf{Совместный} \\
    \hline
    LeNet                  & $1.52$ & $5.21$ \\
    RoBERTa                & $1.76$ & $6.31$ \\
    Внимание (RMSNorm)     & $1.69$--$1.87$ & $5.6$--$6.5$ \\
    SwiGLU / GeGLU FFN     & $1.39$--$1.41$ & $1.72$--$1.76$ \\
    RMSNorm FFN            & $1.47$ & $\mathbf{17.02}$ \\
    DeepNorm               & $\mathbf{0.35}$ & $4.03$ \\
    \hline
\end{tabular}}
\end{table}
